\begin{frame}{정규화의 효과, 숫자로: 0.556 $\to$ 1.000}
  두 클래스가 \textbf{$x$축 하나}로만 분리되는 데이터:
  클래스 0은 $x\approx0.5$, 클래스 1은 $x\approx0.1$,
  $y$축은 두 클래스 모두 $5\pm3$의 잡음.
  $x$의 범위(약 0.6)가 $y$의 범위(약 9)에 비해 매우 작다.
  \vskip0.8em
  \small
  \begin{tabular}{@{}lc@{}}
    \toprule
    & 검증 정확도 ($k=3$) \\
    \midrule
    정규화 \textbf{없이} & 0.556 \\
    정규화 \textbf{후} & \textbf{1.000} \\
    \bottomrule
  \end{tabular}
  \vskip1em
  \begin{block}{해석}
    ``$y$ 잡음으로만 분류하던 모델''이 정규화 후 $x$가 분리 축으로
    정상 작동 $\to$ 정확도 \textbf{거의 두 배}. 정규화 없이
    성능이 40\%나 깎일 수 있는 \textbf{극단적 사례} ---
    센서·금융 데이터 등 실데이터에서 이런 스케일 차이는 흔하다.
  \end{block}
\end{frame}
